% Journal paper set with REVTeX 4.2 for the APS journals and built with
% pdfLaTeX, as the journals and arXiv expect. The class owns the design, so
% the house style is not loaded; the conventions for maths, units, floats and
% references still apply. For another journal, start from its own template.
\documentclass[aps,prb,reprint,superscriptaddress]{revtex4-2}
\usepackage{graphicx}
\usepackage{booktabs}
\usepackage{siunitx}
\usepackage[hidelinks]{hyperref}
% APS style: "Fig. 1" and "Eq. (1)" in running text, spelled out at the start of a sentence.
\usepackage[capitalise]{cleveref}

% ISO 80000-2 sets the differential upright, with a thin space before it.
\newcommand{\dd}{\mathop{}\!\mathrm{d}}

\begin{document}

\title{Measuring the gravitational acceleration with a simple pendulum}
\author{Author Name}
\affiliation{Department of Physics, University Name, City, Country}
\date{\today}

\begin{abstract}
We timed a simple pendulum at nine lengths between \qtylist{0.2;1.0}{\metre} and fitted the small-angle period to the measurements. The fit gives $g = \qty{9.809(4)}{\metre\per\second\squared}$, within one standard error of the local value. The residuals show no trend with length, so the small-angle model describes the pendulum to the timing precision of \qty{1}{\milli\second}.
\end{abstract}

\maketitle

\section{Introduction}

A mass on a light string of length $L$ swings with a period that depends only on $L$ and the gravitational acceleration $g$ when the amplitude is small~\cite{nelson1986}. Timing the period at several lengths therefore measures $g$, and the spread of the measurements about the model tests whether the small-angle approximation holds.

\section{Model and method}

The angle $\theta$ of the string from the vertical obeys
\begin{equation}
  \frac{\dd^2\theta}{\dd t^2} = -\frac{g}{L} \sin\theta ,
  \label{eq:motion}
\end{equation}
which for small $\theta$ gives simple harmonic motion with period
\begin{equation}
  T = 2\pi \sqrt{L/g} .
  \label{eq:period}
\end{equation}
A steel bob on a clamped string was released from rest at \qty{3}{\degree}, where the finite-amplitude correction to \cref{eq:period} is \num{1.7e-4}~\cite{nelson1986}. A light gate timed 20 oscillations at each length, and each period is the mean of five runs with a standard error of \qty{1}{\milli\second}.

\section{Results}

\Cref{fig:period} shows the measured periods with the fit of \cref{eq:period}, and \cref{tab:period} lists them at selected lengths. A least-squares fit of $T = a \sqrt{L}$ gives $g = 4\pi^2 / a^2 = \qty{9.809(4)}{\metre\per\second\squared}$ with $\chi^2 = 6.9$ for 8 degrees of freedom, so the timing error accounts for the scatter~\cite{taylor1997}.

\begin{figure}[tbp]
  \includegraphics[width=\linewidth]{figures/pendulum}
  \caption{Period against length (a), with the fitted small-angle model, and residuals from the fit (b). Error bars are one standard error.}
  \label{fig:period}
\end{figure}

\begin{table}[tbp]
  \caption{Measured and fitted periods at selected lengths.}
  \label{tab:period}
  \begin{tabular}{
    S[table-format=1.1]
    S[table-format=1.4(2)]
    S[table-format=1.4]
  }
    \toprule
    {$L$ (\unit{\metre})} & {$T$ (\unit{\second})} & {$T_\mathrm{fit}$ (\unit{\second})} \\
    \midrule
    0.2 & 0.8983(10) & 0.8972 \\
    0.4 & 1.2692(10) & 1.2688 \\
    0.6 & 1.5548(10) & 1.5540 \\
    0.8 & 1.7937(10) & 1.7943 \\
    1.0 & 2.0057(10) & 2.0061 \\
    \bottomrule
  \end{tabular}
\end{table}

\section{Conclusion}

The residuals scatter about zero with no trend in length, and the fitted $g$ agrees with the local value, so the small-angle model describes the period at every length measured.

\bibliography{refs}

\end{document}
